\section{Estadística y análisis de incertidumbre}

\subsection{Monte Carlo en condiciones iniciales}

Realizar 500 o más simulaciones variando los elementos orbitales iniciales de Apophis dentro de sus barras de error observacionales. Construir el histograma de la distancia mínima durante el encuentro y calcular la probabilidad empírica de que esa distancia sea menor que $30{,}000\,\mathrm{km}$.

\subsection{Mapa de keyholes}

Variar simultáneamente la anomalía media $M_0$ y la excentricidad $e$ sobre una grilla bidimensional y, para cada combinación, calcular la distancia mínima de encuentro en 2029. Representar el resultado con un mapa de calor que muestre las regiones del espacio de parámetros donde el asteroide impactaría.